\begin{helper}[The gradients of the mollified approximating solutions converge in $L^p(B_R)$]
Let $n\ge2$, let $p>1$ and let $r'\in\R$.  Write $p'=p/(p-1)$.  Let $\rho\colon\R^n\to\R$ satisfy
\begin{enumerate}
\item[(H1)] $\rho\ge0$ everywhere and $\int_{\R^n}\rho=1$;
\item[(H2)] $z\mapsto\rho(z)f(z)$ is integrable for every continuous
  $f\colon\R^n\to\R$, and $z\mapsto\rho(z)F(z)$ is integrable for every
  continuous $F\colon\R^n\to\R^n$;
\item[(H3)] $|z|\le1$ for every $z\in\operatorname{supp}\rho$.
\end{enumerate}
Let $A\colon\R^n\to\R^n$ satisfy
\begin{enumerate}
\item[(H4)] $A(0)=0$;
\item[(H5)] $A$ is continuous on $\R^n$.
\end{enumerate}
Let $A_\bullet\colon\R\times\R^n\to\R^n$, written $A_e$, satisfy
\begin{enumerate}
\item[(H6)] $A_e(h)=\int_{\R^n}\rho(z)A(h+\sqrt e\,z)\,dz
  -\int_{\R^n}\rho(z)A(\sqrt e\,z)\,dz$ for all $e\in\R$ and $h\in\R^n$;
\item[(H7)] $A_e(0)=0$ for every $e\in\R$;
\item[(H8)] for every $e$ with $0<e\le1$ there are $\gamma_0',\gamma_1'\in\R$
  with $0<\gamma_0'$, $\gamma_1'\le r'\gamma_0'$ and such that the field
  $(x,z,h)\mapsto A_e(h)$ satisfies \texttt{Q1} and \texttt{Q2} on $\R^n$ with
  the parameters $p,e,\gamma_0',\gamma_1'$
  (\noderef{EllipticStructureQ12}).
\end{enumerate}
Let $x_0\in\R^n$, let $R\in\R$, let $u\colon\R^n\to\R$ and let
$g\colon\R^n\to\R^n$ satisfy
\begin{enumerate}
\item[(H9)] $0<R$;
\item[(H10)] $u\in W^{1,p}(B_R(x_0))$ with weak gradient $g$
  (\noderef{MemW1pOn});
\item[(H11)] $u$ is a weak solution of $\operatorname{div}A(\nabla u)=0$ in
  $B_R(x_0)$ with weak gradient $g$, i.e.\ with the field
  $(x,z,h)\mapsto A(h)$ and vanishing right-hand side
  (\noderef{IsWeakSolutionOn}).
\end{enumerate}
Let $\Gamma\in\R$ satisfy
\begin{enumerate}
\item[(H12)] $\Gamma>0$ and
  $|A(h)|\le\Gamma\bigl(1+|h|^2\bigr)^{\frac{p-1}2}$ for every $h\in\R^n$;
\item[(H13)] $|A_e(h)|\le\Gamma\bigl(1+|h|^2\bigr)^{\frac{p-1}2}$ for every
  $h\in\R^n$ and every $e$ with $0<e\le1$.
\end{enumerate}
Let $c_3\in\R$ satisfy
\begin{enumerate}
\item[(H14)] $c_3>0$ and, for every $e$ with $0<e\le1$ and all $a,b\in\R^n$,
  \[
  c_3\bigl(e+|a|^2+|b|^2\bigr)^{\frac p2-1}|a-b|^2
  \ \le\ \bigl(A_e(a)-A_e(b),\,a-b\bigr).
  \tag{M}
  \]
\end{enumerate}
Finally let $U\colon\mathbb{N}\to(\R^n\to\R)$ and
$G\colon\mathbb{N}\to(\R^n\to\R^n)$ satisfy, for every $m\in\mathbb{N}$ and
with $e_m=1/(m+1)$,
\begin{enumerate}
\item[(H15)] $U_m$ is a weak solution of
  $\operatorname{div}A_{e_m}(\nabla U_m)=0$ in $B_R(x_0)$ with weak gradient
  $G_m$ (\noderef{IsWeakSolutionOn}), and
  $U_m-u\in W^{1,p}_0(B_R(x_0))$ with weak gradient $G_m-g$
  (\noderef{MemW1pZeroOn}).
\end{enumerate}
Then
\[
\int_{B_R(x_0)}|G_m-g|^p\,dx\ \longrightarrow\ 0
\qquad(m\to\infty).
\]
\end{helper}
\begin{proof}
Write $B_R=B_R(x_0)$, $e_m=1/(m+1)$, so that $0<e_m\le1$ for every
$m\in\mathbb{N}$, and
\[
Y_m=\int_{B_R}|G_m|^p\,dx,
\qquad
K=\int_{B_R}|g|^p\,dx,
\qquad
X_m=\Bigl(\int_{B_R}|G_m-g|^p\,dx\Bigr)^{1/p},
\]
\[
I_m=\int_{B_R}\bigl(e_m+|G_m|^2+|g|^2\bigr)^{\frac p2-1}|G_m-g|^2\,dx .
\]
Since $p>1$ we have $p>0$, $p-1>0$ and $p'=p/(p-1)\in(1,\infty)$ with
$(p-1)p'=p$.  The conclusion is $X_m^{\,p}\to0$, and since $t\mapsto t^p$ is
continuous at $0$ with value $0$ it suffices to prove $X_m\to0$.

\emph{Step 0: finiteness, and continuity of the mollified fields.}  By (H15),
$U_m-u\in W^{1,p}(B_R)$ with weak gradient $G_m-g$, so $|G_m-g|\in L^p(B_R)$;
with $|g|\in L^p(B_R)$ from (H10) and the $L^p$ triangle inequality (Minkowski,
Preamble), $|G_m|\in L^p(B_R)$.  Hence $Y_m<\infty$, $K<\infty$ and
$X_m<\infty$, and $|B_R|<\infty$ because $R>0$ is finite.  Moreover, for every
$m$ the field $A_{e_m}$ is continuous on $\R^n$: (H8) at $e=e_m$ supplies an
\noderef{EllipticStructureQ12} bundle for $(x,z,h)\mapsto A_{e_m}(h)$, and the
continuity clause of that predicate, read at any point $x$ and at $z=0$, is
exactly continuity of $A_{e_m}$ as a map $\R^n\to\R^n$.  \emph{This is the only
use made of (H8).}

\emph{Step 1: the two energy identities.}  We claim
\[
\int_{B_R}\bigl(A_{e_m}(G_m),\,G_m-g\bigr)dx=0,
\qquad
\int_{B_R}\bigl(A(g),\,G_m-g\bigr)dx=0 .
\tag{E}
\]
Both are instances of \noderef{WeakSolutionZeroTest}, which extends the
weak-solution identity from the compactly supported test pairs of
\noderef{IsWeakSolutionOn} to the whole of $W^{1,p}_0$ of the domain as soon as
the coefficient lies in $L^{p'}$ there.  That extension is needed because
$U_m-u$ is not of compact support.  We check its hypotheses.

\emph{The coefficients lie in $L^{p'}(B_R)$.}  By (H13) and $(p-1)p'=p$,
\[
\int_{B_R}|A_{e_m}(G_m)|^{p'}dx
\le\Gamma^{p'}\int_{B_R}\bigl(1+|G_m|^2\bigr)^{\frac p2}dx
\le\Gamma^{p'}2^{\frac p2}\bigl(|B_R|+Y_m\bigr)<\infty ,
\tag{1a}
\]
using $(1+t^2)^{p/2}\le2^{p/2}(1+t^p)$ for $t\ge0$ and $Y_m<\infty$,
$|B_R|<\infty$ from Step 0; the integrand is measurable because $A_{e_m}$ is
continuous by Step 0 and $G_m$ is measurable.  The same computation with (H12)
in place of (H13) and $K$ in place of $Y_m$ gives
$|A(g)|\in L^{p'}(B_R)$, and the one with (H13) at the argument $g$ gives
$|A_{e_m}(g)|\in L^{p'}(B_R)$, again with $K$ in place of $Y_m$; the latter is
used in Step 3.

\emph{The test pair is admissible.}  By (H15), $U_m-u\in W^{1,p}_0(B_R)$ with
weak gradient $G_m-g$ (\noderef{MemW1pZeroOn}).

Now apply \noderef{WeakSolutionZeroTest} twice, with $1\le n$ (from $n\ge2$),
$1<p$, the open set $B_R$ and the test pair $(U_m-u,\,G_m-g)$: first to the
field $(x,z,h)\mapsto A_{e_m}(h)$ and the solution $(U_m,G_m)$ of (H15), whose
coefficient membership is (1a); then to the field $(x,z,h)\mapsto A(h)$ and the
solution $(u,g)$ of (H11), whose coefficient membership is the second case of
(1a).  Both weak-solution hypotheses have vanishing right-hand side, so the
cited identity is exactly $\int_{B_R}\sum_i(A_\bullet)_i(G_m-g)_i\,dx=0$ in each
case, which is (E).

\emph{Step 2: a bound on $Y_m$ uniform in $m$.}  The first identity in (E)
reads $\int_{B_R}(A_{e_m}(G_m),G_m)=\int_{B_R}(A_{e_m}(G_m),g)$.  For the
left-hand side use (M) with $a=G_m(x)$, $b=0$ and $A_{e_m}(0)=0$ from (H7):
\[
\bigl(A_{e_m}(G_m(x)),G_m(x)\bigr)
\ \ge\ c_3\bigl(e_m+|G_m(x)|^2\bigr)^{\frac p2-1}|G_m(x)|^2 .
\]
If $p\ge2$ the exponent is nonnegative, so
$(e+|h|^2)^{\frac p2-1}\ge(|h|^2)^{\frac p2-1}=|h|^{p-2}$ and the right-hand
side is at least $c_3|h|^p$ with $h=G_m(x)$.  If $1<p<2$ the exponent is
negative; then at points where $|h|^2\ge e$ we have $e+|h|^2\le2|h|^2$ and hence
$(e+|h|^2)^{\frac p2-1}\ge2^{\frac p2-1}|h|^{p-2}$, so the right-hand side is at
least $c_32^{\frac p2-1}|h|^p$ there, while at the remaining points it is
nonnegative and $|h|^p<e^{p/2}$.  With $\kappa_p=\min\{1,2^{\frac p2-1}\}>0$,
integrating and using monotonicity of the integral (Preamble),
\[
\int_{B_R}\bigl(A_{e_m}(G_m),G_m\bigr)dx
\ \ge\ c_3\kappa_p\Bigl(Y_m-e_m^{\frac p2}|B_R|\Bigr).
\]
For the right-hand side, Hölder's inequality and the computation of Step 1 give
\[
\Bigl|\int_{B_R}\bigl(A_{e_m}(G_m),g\bigr)dx\Bigr|
\le\Gamma2^{\frac p{2p'}}\bigl(|B_R|+Y_m\bigr)^{1/p'}K^{1/p}
=:b_0\bigl(|B_R|+Y_m\bigr)^{\frac{p-1}p},
\]
with $b_0=\Gamma2^{\frac p{2p'}}K^{1/p}$ independent of $m$.  Putting
$S=|B_R|+Y_m$ and using $e_m\le1$,
\[
c_3\kappa_pS\ \le\ c_3\kappa_p\bigl(|B_R|+e_m^{\frac p2}|B_R|\bigr)+b_0S^{\theta}
\ \le\ 2c_3\kappa_p|B_R|+b_0S^{\theta},
\qquad \theta=\frac{p-1}p\in(0,1),
\]
that is, $S\le a_0+b_1S^{\theta}$ with $a_0=2|B_R|$ and
$b_1=b_0/(c_3\kappa_p)$.  If $S>2a_0$ then $S/2\le b_1S^{\theta}$, so
$S^{1-\theta}\le2b_1$ and $S\le(2b_1)^{1/(1-\theta)}$.  Hence
\[
|B_R|+Y_m\ \le\ E:=\max\bigl\{2a_0,\ (2b_1)^{1/(1-\theta)}\bigr\}
\qquad(m\in\mathbb{N}),
\tag{2a}
\]
a bound independent of $m$.  Consequently, using
$(s+t+w)^{p/2}\le3^{p/2}(s^{p/2}+t^{p/2}+w^{p/2})$ for $s,t,w\ge0$ and
$e_m\le1$,
\[
\int_{B_R}\bigl(e_m+|G_m|^2+|g|^2\bigr)^{\frac p2}dx
\ \le\ 3^{\frac p2}\bigl(|B_R|+Y_m+K\bigr)\ \le\ E_1:=3^{\frac p2}(E+K)
\qquad(m\in\mathbb{N}).
\tag{2b}
\]

\emph{Step 3: the difference estimate.}  Subtracting the two identities of (E),
\[
\int_{B_R}\bigl(A_{e_m}(G_m)-A_{e_m}(g),\,G_m-g\bigr)dx
=\int_{B_R}\bigl(A(g)-A_{e_m}(g),\,G_m-g\bigr)dx ,
\]
because the first identity makes $\int(A_{e_m}(G_m),G_m-g)=0$ and the second
makes $\int(A(g),G_m-g)=0$, so both sides equal
$-\int(A_{e_m}(g),G_m-g)$; all four integrals are finite by Step 1 and Hölder,
the bound of Step 1 applying to $A_{e_m}(g)$ as well, with $Y_m$ replaced by
$K$.  By (M) with $a=G_m(x)$, $b=g(x)$ the left-hand side is at least $c_3I_m$,
and by Hölder's inequality the right-hand side is at most $\delta_mX_m$, where
\[
\delta_m=\Bigl(\int_{B_R}\bigl|A(g)-A_{e_m}(g)\bigr|^{p'}dx\Bigr)^{1/p'} .
\]
Hence
\[
c_3I_m\ \le\ \delta_mX_m\qquad(m\in\mathbb{N}).
\tag{3}
\]

\emph{Step 4: $\delta_m\to0$.}  First, $A_e(h)\to A(h)$ for every fixed $h$ as
$e\downarrow0$.  Indeed $A$ is continuous on $\R^n$ by (H5), so
$\sup_{|y-h|\le\sqrt e}|A(y)-A(h)|\to0$ as $e\downarrow0$ by continuity at $h$;
since $|z|\le1$ on $\operatorname{supp}\rho$ by (H3), $\rho\ge0$ and
$\int\rho=1$ by (H1), and the integrands are integrable by (H2) --- the maps
$z\mapsto A(h+\sqrt e\,z)$ and $z\mapsto A(\sqrt e\,z)$ being continuous ---
this gives
\[
\int_{\R^n}\rho(z)A(h+\sqrt e\,z)\,dz\longrightarrow A(h)
\qquad\text{and}\qquad
\int_{\R^n}\rho(z)A(\sqrt e\,z)\,dz\longrightarrow A(0)=0
\]
by (H4), the error in each case being at most
$\int\rho(z)\sup_{|y-h|\le\sqrt e}|A(y)-A(h)|\,dz$ over
$\operatorname{supp}\rho$.  Subtracting and recalling (H6) gives the claim.
Second, by (H12) and (H13),
\[
\bigl|A(g(x))-A_{e}(g(x))\bigr|^{p'}
\ \le\ (2\Gamma)^{p'}\bigl(1+|g(x)|^2\bigr)^{\frac p2}
\qquad(0<e\le1),
\]
and the majorant is integrable on $B_R$, being at most
$2^{p/2}(2\Gamma)^{p'}(1+|g|^p)$ with $|B_R|<\infty$ and $|g|\in L^p(B_R)$.  The
integrands are measurable, $A$ and each $A_{e_m}$ being continuous by (H5) and
Step 0.  Since $e_m=1/(m+1)\downarrow0$, dominated convergence (Preamble)
applied to the sequence $x\mapsto|A(g(x))-A_{e_m}(g(x))|^{p'}$ --- which
converges pointwise to $0$ by the first part --- gives $\delta_m\to0$.

\emph{Step 5: $X_m\le\Psi(\delta_m)$ with $\Psi(\delta)\to0$ as $\delta\to0$.}
Fix $m$.  If $X_m=0$ the bounds below hold trivially, so assume $X_m>0$.

\emph{Case $p\ge2$.}  Pointwise $|a-b|^2\le2|a|^2+2|b|^2\le2(e+|a|^2+|b|^2)$ for
$e\ge0$, so $e+|a|^2+|b|^2\ge\frac12|a-b|^2$ and, the exponent $\frac p2-1$
being nonnegative,
$\bigl(e+|a|^2+|b|^2\bigr)^{\frac p2-1}\ge2^{1-\frac p2}|a-b|^{p-2}$.  Hence
$I_m\ge2^{1-\frac p2}X_m^{\,p}$, and (3) gives
$c_32^{1-\frac p2}X_m^{\,p}\le\delta_mX_m$, i.e.
\[
X_m\ \le\ \Bigl(\frac{2^{\frac p2-1}}{c_3}\,\delta_m\Bigr)^{\frac1{p-1}} .
\]

\emph{Case $1<p<2$.}  Since $e_m>0$, the quantity
$\Theta(x)=\bigl(e_m+|G_m(x)|^2+|g(x)|^2\bigr)^{\frac p2-1}$ is finite and
strictly positive, so we may write
$|G_m-g|^p=\bigl(\Theta|G_m-g|^2\bigr)^{\frac p2}\Theta^{-\frac p2}$ pointwise.
Hölder's inequality with the conjugate exponents $\frac2p$ and $\frac2{2-p}$
(Preamble) gives
\[
X_m^{\,p}=\int_{B_R}\bigl(\Theta|G_m-g|^2\bigr)^{\frac p2}\Theta^{-\frac p2}dx
\ \le\ I_m^{\frac p2}
\Bigl(\int_{B_R}\Theta^{-\frac p{2-p}}dx\Bigr)^{\frac{2-p}2} .
\]
Now $-\frac p{2-p}\bigl(\frac p2-1\bigr)=\frac p{2-p}\cdot\frac{2-p}2=\frac p2$,
so $\Theta^{-\frac p{2-p}}=\bigl(e_m+|G_m|^2+|g|^2\bigr)^{\frac p2}$ and the
last integral is at most $E_1$ by (2b).  With (3),
\[
X_m^{\,p}\ \le\ \Bigl(\frac{\delta_mX_m}{c_3}\Bigr)^{\frac p2}E_1^{\frac{2-p}2},
\qquad\text{hence}\qquad
X_m^{\frac p2}\le\Bigl(\frac{\delta_m}{c_3}\Bigr)^{\frac p2}E_1^{\frac{2-p}2},
\qquad\text{i.e.}\qquad
X_m\ \le\ \frac{\delta_m}{c_3}\,E_1^{\frac{2-p}p} .
\]

\emph{Step 6: conclusion.}  In both cases $X_m$ is bounded by a continuous
increasing function of $\delta_m$ that vanishes at $\delta_m=0$, and the bound
is uniform in $m$ --- in the case $p\ge2$ it is
$\delta\mapsto(2^{p/2-1}\delta/c_3)^{1/(p-1)}$ and in the case $1<p<2$ it is
$\delta\mapsto\delta E_1^{(2-p)/p}/c_3$, the constants $c_3$ and $E_1$ not
depending on $m$ by (H14) and (2a).  Since $\delta_m\to0$ by Step 4, this gives
$X_m\to0$, and therefore
\[
\int_{B_R}|G_m-g|^p\,dx=X_m^{\,p}\longrightarrow0 ,
\]
which is the assertion.
\end{proof}
